\section{Многообразия простых особенностей.}
Рассмотрим простую особенность типа \(A_{\mu}\) заданную многочленом \(f_{A_{\mu}} = \frac{1}{\mu+1}x^{\mu +1}\). Ассоциированное с ним многообразие Дубровина-Фробениуса определим так:\par 
рассмотрим многочлен \[\Lambda = f_{A_{\mu}} + \sum_{0 \leq k < \mu} s_k x^k\] и фактор \[\mathbb{C}[x, s_0, s_1, \dots, s_\mu]\Big/(\partial_x \Lambda)\]  

Для многообразия \(M = \mathbb{C}^{\mu} = (s_0, s_1, \dots, s_{\mu-1})\), касательным векторам \(\tang[s^m]\) сопоставляются соответствующие частные производные многочлена \(\Lambda\), \(\tang[s^m] \mapsto	\dfrac{\partial\Lambda}{\partial s^m}\). Умножение задается умножением в кольце многочленов \(\mathbb{C}[x, s_0, \dots, s_{\mu-1}]/\partial_x \Lambda\).
\[\tang[s^a]\circ\tang[s^b] = \dfrac{\partial\Lambda}{\partial s^a} \cdot \dfrac{\partial\Lambda}{\partial s^b}, \operatorname{modulo} \partial_x\Lambda\]

\subsection{Градуировка многообразия \(M_{A_\mu}\)}
В случае \(A_\mu\), \(\Lambda(x,s)=\dfrac{x^{\mu+1}}{\mu+1}+\sum_{i=0}^{\mu-1}s_i x^i.\)
Касательная алгебра имеет вид\[\mathbb{C}[x]\Big/
\left(\frac{\partial\Lambda}{\partial x}\right)\]
\[
e_i=\frac{\partial}{\partial s_i}
\longmapsto
\frac{\partial\Lambda}{\partial s_i}
=
x^i.
\]
Следовательно, на стороне аналитического спектра \(y_i|_{L_M}=x^i\)
Поэтому
\[\operatorname{wt}_{T_M}(e_i)=\operatorname{wt}_{\mathcal O_{L_M}}(x^i).\]
А значит
\begin{equation} \label{grad}
	\begin{gathered}
		\operatorname{wt}(x)=\frac{1}{\mu+1}\Rightarrow\\
		\Rightarrow \operatorname{wt}(e_i)=\frac{i}{\mu+1}\Rightarrow\\
		\Rightarrow \operatorname{wt}(s_i)=1-\frac{i}{\mu+1}.
	\end{gathered}
\end{equation}
Сформулируем и докажем следующую лемму.
\begin{lem}\label{lem1}
	Пусть дано некоторое многообразие Дубровина-Фробениуса \((M, \circ)\) размерности \(\mu\) с градуировкой \(A_\mu\).
	\begin{enumerate}
		\item Существует семейство фактор алгебр, зависящих от \(s\), таких что 
		\[T_pM \simeq B = \mathbb{C}[x]/(x^{\mu} + \sum \alpha_k x^k), \text{где } \alpha_i \in \mathbb{C}[s_0, \dots, s_{\mu-1}]\]
		\item Для любого изоморфизма \(\mathbb{Q}\)-градуированных алгебр \(\varphi:T_0M \rightarrow B|_0\), выполнено \(\varphi\left(\dfrac{\partial}{\partial s^\alpha}\right) = c_\alpha x^i\)
	\end{enumerate}
\end{lem}
\begin{proof}
	\(T_pM\) порождено касательными векторами \(\langle \partial/\partial s^0, \dots, \partial/\partial s^{\mu-1} \rangle\) c соотношениями \[\tang[s^\alpha] \circ \tang[s^\beta] = c^\gamma_{\alpha\beta}\tang[s^{\gamma}]\]
	Построим в \(T_pM\) новый базис \[\{e_0, e_1, \dots, e_1^{\mu-1}\}\] индуктивно. \par
	\paragraph{База индукции: $k=0,1$}
	\begin{itemize}
		\item В качестве \(e_0\) выберем единичное поле умножения \(\circ\). \(e_0 = e\).
		\item В \(T_pM\) градуировка выделяет одномерное подпространство векторов наименьшего возможного веса \(\dfrac{1}{3}\). Выберем в нем любой ненулевой вектор в качестве \(e_1\). 
		
	\end{itemize}
	Коэффициенты полиномиальны (константы).
	Для изоморфизма $\varphi$:
	\begin{itemize}
		\item $\varphi(e_0)=\varphi(1)=1=c_0x^0$ $(c_0=1\neq 0)$;
		\item $\varphi(e_1)=c_1x; \quad c_1\in \mathbb{C}\setminus\{0\})$, так как
		\(\operatorname{wt}(e_1)=\frac{1}{n+1}=\operatorname{wt}(x),\)
		а элементов с меньшим весом или хотя бы таким же нет.
	\end{itemize}
	База доказана.
	\paragraph{Переход}
	Пусть \(\forall j \leq k\) такой базис есть, тогда 
	\begin{gather*}
		\operatorname{wt}(e_k\cdot e_1) = \dfrac{k}{\mu + 1} + \dfrac{1}{\mu + 1} = \dfrac{k + 1}{\mu + 1}\\
		\operatorname{wt}(c^t_{k1}) = \dfrac{k+1}{\mu+1} - \dfrac{t}{\mu +1} = \dfrac{k+1-t}{\mu + 1} \text{ или } c^t_{k1} = 0
	\end{gather*}
	Выразим \(e_{k+1}\): 
	\begin{align*}
		e_k \cdot e_1 = c_{k1}^{k+1} e_{k+1} + \sum^k_{t = 0} c^t_{k1} e_t &\Rightarrow\\
		\Rightarrow & \hspace{3pt}e_{k+1} = \dfrac{1}{c_{k1}^{k+1}} \left( e_k\cdot e_1  - \sum^k_{t = 0} c^t_{k1} e_t \right)
	\end{align*}
	Заметим,что 
	\(c^t_{k1}(0) = 0\) если \(t \neq k+1\), а \(0 \neq c^{k+1}_{k1} = const \in \mathbb{C}\)
	Это верно, так как 
	\begin{itemize}
		\item \(t > k+1 \Rightarrow\operatorname{wt}(c^t_{k1}) < 0 \Rightarrow  c^t_{k1} = 0\)
		\item \(t < k+1 \Rightarrow\operatorname{wt}(c^t_{k1}) > 0 \Rightarrow c^t_{k1}(0) = 0\), так как зануляются свободные члены.
		\item \(t = k+1 \Rightarrow\operatorname{wt}(c^t_{k1}) = 0 \Rightarrow c^t_{k1} = const\)\\
		Пользуясь предыдущими пунктами при ограничении в ноль получаем \[e_ke_1 = c_{k1}^{k+1}e_{k+1}\] так как все остальное обращается в ноль. Для изоморфизма \(\varphi\) по предположению индукции получаем 
		\[c_{k1}^{k+1}\varphi(e_{k+1}) = \varphi(c_{k1}^{k+1}e_{k+1}) =\varphi(e_ke_1) = \varphi(e_k)\varphi(e_1) = c_kc_1x^{k+1} \neq 0 \Rightarrow c_{k1}^{k+1} \neq 0\]
	\end{itemize}
	Из этого замечания выражение выше корректно. Значит, пользуясь предположением индукции получаем, что \(e_{k+1}\) выражается через \(e_0, e_1, \dots, e_1^{k+1}\).\\
	Рассмотрим исходное структурное равенство в новом базисе:
	\[e_{\mu-1} e_1 = \sum_{k = 0}^{\mu -1} c^k_{n-1,1} e_k\]
	Подставляя выражения для \(e_{\mu -1}, e_k\) через степени \(e_1\) и приводя подобные получаем: \[e_1^{\mu}- \sum^{\mu -1}_{k = 0}\alpha_k(s) = 0, \quad \alpha_k(s) \in \mathbb{C}[s_0, \dots, s_{\mu-1}]\]
	Т.е. имеет место изоморфизм
	\[T_pM \simeq \mathbb{C}[x]/(x^{\mu} + \sum \alpha_k x^k), \text{где } \alpha_i \in \mathbb{C}[s_0, \dots, s_{\mu-1}],\] устанавливающий взаимно-однозначное соответствие между базисами \(\{1, x, \dots, x^{\mu-1}\}\) и \(\{e_0, e_1, \dots, e^{\mu-1}_1\}\).
\end{proof}
Итак, имеем структурное соотношение 
\begin{equation}
	e_1^{\circ(\mu)}
	=
	\sum_{k=0}^{\mu-1}\alpha_k(s)e_k \quad \text{где } \quad \alpha_k(0)=0.
	\label{eq:fundamental_relation} 
\end{equation}
